% Transcription — fonds Alexandre Grothendieck, shelfmark 158, batch 4 (pages 61-80).
% Established from the facsimile archives/batches/158/batch-04.pdf (a mirror of
% https://grothendieck.umontpellier.fr/158.pdf), per the skill
% transcribe-grothendieck. The batch file carries no cover sheet of its own:
% PDF page N = archivists' page 60+N, checked against the pencilled numbers
% bottom-left (PDF p. 1 bears « 61 », PDF p. 20 bears « 80 »).
% Pages 63, 64 and 65 are one administrative circular, page 80 a printed
% seminar announcement: no mathematics — absent.
% Page 68 carries mathematics on its upper half only; the lower half is an
% address list.
% The run is his own pp. 2-11, so his p. 1 falls in batch 3.
% Pass: Opus 5 (claude-opus-5), 2026-09-20 — first pass, unchecked against the pages by a human.
\documentclass[11pt,a4paper]{article}
% Le préambule commun à toutes les transcriptions du fonds Grothendieck.
%
% Il tient en une idée : **l'incertitude se compose, elle ne se résout pas.**
% Une transcription qui lisse un mot douteux a détruit la seule chose pour
% laquelle on la faisait — un lecteur doit pouvoir savoir, à chaque endroit,
% ce qui était sur le feuillet et ce qui a été ajouté. Les six macros qui
% suivent sont donc l'appareil critique, et non de la décoration.
%
% Les mêmes commandes sont reconnues par `scripts/render.mjs`, qui produit la
% vue de lecture du site. Une macro ajoutée ici doit l'être aussi là-bas,
% faute de quoi le rendu échouera bruyamment — ce qui est voulu : un rendu
% silencieusement faux est pire qu'un rendu absent.

\ProvidesPackage{grothendieck}[2026/08/08 Transcriptions du fonds Grothendieck]

\usepackage{amsmath,amssymb,amsthm}
\usepackage{mathrsfs}
\usepackage{stmaryrd}
% Les diagrammes commutatifs sont partout dans ce fonds : c'est la langue
% même de la géométrie algébrique de Grothendieck, et une transcription qui
% les rendrait en prose ne transcrirait rien.
\usepackage{tikz-cd}
% Français par défaut — c'est la langue des feuillets — mais l'anglais est
% chargé aussi : la traduction et le résumé sont des documents anglais, et la
% typographie française (espaces avant les deux-points) y serait une faute.
% Ces documents basculent par \selectlanguage{english} en tête de corps.
\usepackage[english,french]{babel}
\usepackage{fontspec}
\usepackage{geometry}
\geometry{a4paper,margin=25mm,marginparwidth=28mm}
\usepackage{marginnote}
\usepackage{xcolor}
% ulem, not soul. `soul`'s \st and \ul are built on a character-by-character
% scan that XeLaTeX's font machinery breaks: the document fails with an
% undefined control sequence rather than degrading. ulem does the same two jobs
% and is Unicode-engine safe. `normalem` stops it hijacking \emph, which the
% transcriptions use for Grothendieck's own emphasis.
\usepackage[normalem]{ulem}
\usepackage{hyperref}
\hypersetup{colorlinks=true,linkcolor=black,urlcolor=black,pdfborder={0 0 0}}

% --- Métadonnées du lot -----------------------------------------------------
% Reprises telles quelles par le script de rendu, qui les affiche en tête de
% la vue de lecture. Elles fixent ce que le fichier prétend couvrir : sans
% elles, une transcription flotte sans point d'attache dans un fonds de seize
% mille feuillets.

\newcommand{\folder}[1]{\gdef\grothFolder{#1}}
\newcommand{\batch}[1]{\gdef\grothBatch{#1}}
\newcommand{\leaves}[2]{\gdef\grothFirst{#1}\gdef\grothLast{#2}}
\newcommand{\foldertitle}[1]{\gdef\grothFoldertitle{#1}}
\gdef\grothFolder{?}\gdef\grothBatch{?}\gdef\grothFirst{?}\gdef\grothLast{?}\gdef\grothFoldertitle{}

% --- Le résumé --------------------------------------------------------------
%
% Il ouvre la lecture modernisée, sous le titre « Résumé », et il est composé
% en petit. Ce n'est pas un ornement : c'est le seul passage du document qui
% ne soit pas des mathématiques, et un lecteur doit pouvoir le reconnaître
% avant de l'avoir lu — puis le sauter s'il connaît déjà le sujet, ou s'y
% arrêter s'il ne le connaît pas. Le filet qui le suit marque la reprise du
% corps.

\newenvironment{resume}
  {\small\setlength{\parskip}{0.45em}}
  {\par\medskip\hrule\medskip}

% --- Mots-clés ---------------------------------------------------------------
%
% En anglais, en clôture du résumé de la lecture modernisée : le vocabulaire
% moderne sous lequel chercher ce que les feuillets construisent. Le manifeste
% du site (`npm run manifest`) les extrait pour étiqueter la cote entière —
% cette ligne est la source unique des tags, il n'y a pas de fichier à côté.
\newcommand{\keywords}[1]{\par\smallskip\noindent
  {\footnotesize\textsc{Keywords} --- \textit{#1}}\par}

% --- L'appareil critique ----------------------------------------------------

% Le numéro de feuillet, en marge. C'est l'ancre qui permet de régler un
% désaccord sur une formule en regardant *un* feuillet plutôt qu'une chemise
% de six cents. Le site s'en sert aussi pour tourner les pages du fac-similé
% quand on fait défiler la transcription.
\newcommand{\leaf}[1]{%
  \marginnote{\footnotesize\color{gray!70}\textsf{#1}}%
  \label{leaf:#1}\ignorespaces}

% Illisible : on ne devine pas. Un mot inventé qui se lit comme les autres est
% le pire résultat possible.
\newcommand{\ill}{\textcolor{red!60!black}{[\,\ldots\,]}}

% Lecture douteuse : le mot est donné, et signalé comme incertain.
\newcommand{\uncertain}[1]{\uline{#1}}

% Ajout de l'éditeur — une lettre manquante, un mot sous-entendu.
\newcommand{\add}[1]{\textcolor{blue!50!black}{[#1]}}

% Biffé par Grothendieck lui-même. Ce qu'il a rayé fait partie du document :
% une rature dit souvent où la pensée a changé de direction.
\newcommand{\struck}[1]{\sout{#1}}

% Note du transcripteur, sur l'état matériel du feuillet.
\newcommand{\note}[1]{\par\smallskip{\small\color{gray!60!black}\textit{[#1]}}\par\smallskip}

% Note marginale de Grothendieck, distincte du corps du texte.
\newcommand{\marginal}[1]{\par\smallskip\noindent\hspace{1em}%
  {\small\color{gray!60!black}#1}\par\smallskip}

% --- Titre ------------------------------------------------------------------

\AtBeginDocument{%
  \begin{center}
    {\large\bfseries Fonds Alexandre Grothendieck --- cote n\textsuperscript{o}~\grothFolder}\\[2pt]
    {\itshape\grothFoldertitle}\\[4pt]
    {\small feuillets \grothFirst--\grothLast\ (lot \grothBatch)}\\[2pt]
    {\footnotesize\color{gray!60!black}Transcription établie d'après le fac-similé numérisé
     par l'Université de Montpellier.\\
     Les lectures incertaines sont soulignées, les illisibles notées [\,\ldots\,],
     les ajouts entre crochets.}
  \end{center}
  \vspace{1em}}

\endinput


\folder{158}
\batch{4}
\pages{61}{80}
\dating{[vers 1990-1991]}
\watermark{Édition de démonstration}
\foldertitle{[Autour des Dérivateurs] : notes manuscrites (s.d.)}

\begin{document}

\section*{Localisateurs, 0-connexité et suspensions de catégories}
\note{titre de l'éditeur : les feuillets n'en portent aucun. Le lot reprend
une suite paginée par lui de 2 à 11 (feuillets 61, 62, 66, 68, 70, 72, 74, 76,
78, 79) ; le feuillet 1 de sa suite tombe dans le lot précédent}

\note{Notations tenues pour tout le lot, telles qu'elles sont tracées.
$X^{\wedge}$ et $X^{\vee}$ : les deux catégories de foncteurs à valeurs
ensembles sur $X$ (préfaisceaux, resp. covariants) ; $X_{/F}$ et $X_{/a}$ :
catégorie au-dessus ; ${}_{x}\backslash X$ : catégorie au-dessous, l'indice
porté à gauche d'une barre oblique inverse (feuillets 74 et 77, où les deux
sens figurent dans la même formule) ; $A/\!/B$ : signe à deux jambages
obliques, porté tantôt en préfixe ($\beta/\!/Z$, $b_{0}/\!/X_{0}$, feuillets
62 et 66), tantôt en indice postfixe ($\tilde X_{/\!/\tilde\xi}$, feuillets 76
et 79). $W$ : un localisateur, avec les axiomes qu'il numérote Loc(6), Loc(8) ;
$W_{2}$, $W_{\infty}$ : deux localisateurs particuliers. $\Sigma^{n}$ : la
suspension itérée, $e$ l'objet (ou la catégorie) final, $E$ l'ensemble à deux
éléments $\{\beta,\gamma\}$. J'écris $/\!/$ et $\backslash$ comme il les
trace, sans les uniformiser}

\page{61}
\note{il numérote ce feuillet « 2 »}

\begin{tikzcd}
x \arrow[r, bend left=14] \arrow[r, bend right=14] & a
\end{tikzcd}

\[
  \pi_{0}(S_{/x}) \rightrightarrows \pi_{0}(S_{/a})
\]

\note{suivent deux triangles esquissés, dont les sommets ne se lisent pas}

\[
  \pi_{0}(E_{/x}) \simeq
\]
\ill{}

\begin{tikzcd}
\mathrm{Hom}(a,x)\times E \arrow[d] \\
\mathrm{Hom}(a,a)\times E
\end{tikzcd}

\note{le second argument du $\mathrm{Hom}$ inférieur est douteux : je lis $a$,
mais le trait souffrirait un autre sommet. La flèche verticale est doublée
d'un second trait, sinueux, sans tête}

Soit $X$ dans $\mathrm{Cat}$ \ill{} (quelconque) \ill{}, $B$, $C$, $\beta$,
$\gamma$ tels que \ill{}$(\beta,\beta)$, $\mathrm{Hom}(\beta,\gamma)$,
$\mathrm{Hom}($\ill{}$)$ \ill{}

\[
  B\times C = \emptyset \quad\text{dans } X^{\wedge},
\]
i.e. \ill{} n'existe pas de $x\in\mathrm{Ob}\,X$ \ill{} tel que
$\mathrm{Hom}(x,B)\neq\emptyset$ et $\mathrm{Hom}(x,C)\neq\emptyset$,
et
\[
  \beta\times\gamma = \emptyset \quad\text{dans } X^{\vee},
\]
i.e. n'existe pas de $y\in\mathrm{Ob}\,X$ \ill{} $\mathrm{Hom}(\beta,y)\neq
\emptyset$, $\mathrm{Hom}(\gamma,y)\neq\emptyset$. (\ill{})
\note{le signe entre $\beta$ et $\gamma$ est tracé à deux jambages, et non
comme le $\times$ de la ligne précédente ; c'est la glose qui suit qui donne
le sens — le produit des deux foncteurs représentables dans $X^{\vee}$}

Alors $X$ n'est pas $1$-connexe.
\note{les huit dernières lignes du feuillet sont portées en biais et
d'une écriture très rapide ; c'est le passage le plus incertain du lot}

\page{62}
\note{il numérote ce feuillet « 3 »}

\begin{tikzcd}
& b \\
a \arrow[ru] \arrow[r, "u"] \arrow[rd] & x \\
& c
\end{tikzcd}

\note{les sommets $b$ et $c$ sont hachurés}

\[
  \mathrm{Hom}(a,x) = \emptyset
\]
\[
  (u,\beta) \in \mathrm{Hom}(a,x)\times E
\]
\ill{} mais alors $\pi_{0}(S_{/x}) = \emptyset$, i.e. \ill{}

\note{un cadre, écrit en travers du feuillet, porte une formule dont je ne
lis avec quelque assurance que $e\backslash X = $ \struck{\ill{}} $= Z$, et
plus bas $T$, $\beta/\!/Z$ et $C/\!/Z$}

\[
  \pi_{0}(S_{/x}) \simeq I(x) \amalg \bigl(J(x)\times E\bigr)
\]
\note{un membre biffé, illisible, précède $I(x)$}
\[
  I(x) = \mathrm{Im}\bigl(\mathrm{Hom}(\beta,x)\amalg\mathrm{Hom}(\gamma,x)
  \longrightarrow \mathrm{Hom}(a,x)\bigr)
\]
\note{le nom de l'opération se lit « Im » ou « Inv »}
\[
  J(x) = \mathrm{Hom}(a,x)\setminus I(x)
\]
\[
  \pi_{0}(S_{/x}) = \emptyset \quad\text{si } \mathrm{Hom}(a,x) = \emptyset
\]

Soit $\xi = \mathrm{cl}(u,\alpha) \in$ \ill{}, $u \in \mathrm{Hom}(a,x)$,
$\alpha \in E = \{\beta,\gamma\}$. Alors \ill{} $\alpha \longrightarrow \xi$,
i.e. $\xi \in \alpha/\!/Z$.

\begin{tikzcd}[column sep=small, row sep=small]
& Z = e\backslash X \sim S^{1} \arrow[ld, "\text{asph.}"'] \arrow[rd, "\text{asph.}"] & \\
e \arrow[rd] & & X \arrow[ld] \\
& X' \sim S^{2} &
\end{tikzcd}

\[
  X' \sim \Sigma^{1} Z
\]

\page{66}
\note{il numérote ce feuillet « 4 »}

\note{le haut du feuillet porte un losange dont le sommet se lit
$(b\amalg c)$ \ill{} et dont les deux côtés descendants, étiquetés
« asph.\,? », aboutissent à $e\backslash X = Z$ ; je ne lis pas les deux
sommets intermédiaires, l'un étant biffé}

\begin{tikzcd}[column sep=small, row sep=small]
& e\backslash X = Z \arrow[ld, "\text{asph.}"'] \arrow[rd] & \\
e \arrow[rd] & & e\backslash X \arrow[ld, "\text{asph.}"] \\
& X' &
\end{tikzcd}
\note{le sommet de droite est biffé sur le feuillet}

\[
  \beta/\!/Z \simeq {}_{a}\backslash X
\]

\begin{tikzcd}[column sep=small, row sep=small, nodes={font=\scriptsize}]
& (b_{0}\amalg c_{0})/\!/X_{0} \arrow[ld] \arrow[rd] & \\
b_{0}/\!/X_{0} \arrow[rd, "W\text{-cocart.}"'] & & c_{0}/\!/X_{0} \arrow[ld] \\
& T \arrow[ld] \arrow[rd] & \\
X_{0} = {}_{a}\backslash X \arrow[rd, hook, "W\text{-cocart.}"'] & & {}_{a}\backslash X = X_{0} \arrow[ld] \\
& e\backslash X = \tilde X \arrow[ld] \arrow[rd, "\text{coét.}"] & \\
e \arrow[rd, "W\text{-cocart.\,!}"'] & & X \arrow[ld] \\
& X' &
\end{tikzcd}
\note{le premier sommet du losange supérieur porte, entre parenthèses, un
signe $\amalg$ ; le sommet $T$ est suivi d'un membre de droite entièrement
biffé, réécrit à la ligne suivante}

\[
  T = b_{0}/\!/X_{0} \ \cup\ c_{0}/\!/X_{0}
\]
\[
  X_{0} \simeq \beta/\!/\tilde X \simeq \gamma/\!/\tilde X
\]
\[
  \tilde X \sim \Sigma^{1}T \qquad X' \sim \Sigma^{2}T \qquad
  X' \simeq \Sigma^{1}\tilde X
\]
\note{ici $\tilde X$ est \emph{défini} par $e\backslash X$ ; au feuillet 76 il
écrit au contraire $\tilde X = X_{/e}$. Je transcris chaque fois ce qui est
tracé}

\page{67}

\[
  \Sigma^{1}X
\]
\note{un dessin : deux cônes opposés par leur base, la base étant un disque
hachuré — la suspension au sens usuel ; à côté, son avatar combinatoire,
un losange de quatre sommets, légendé $(X,e)$ avec une flèche vers $e$}

\begin{tikzcd}[column sep=small, row sep=small]
& X_{0} \arrow[ld] & \\
X_{1} \arrow[rd] & & X_{2} \arrow[ld] \arrow[lu] \\
& X_{3} &
\end{tikzcd}

\[
  X\times Y \qquad X \longrightarrow X\times I \qquad x_{0}\times I
\]
\note{feuillet de croquis : rien n'y est rédigé}

\page{68}
\note{il numérote ce feuillet « 5 ». La moitié inférieure du feuillet ne
relève pas des notes}

\note{en haut, un circuit de flèches où se lisent $X_{/F} \simeq X_{/e}$,
$e$, $X$, $X'$, et les mentions « cocoasph. », « ouv. », « ferm. » ;
plusieurs sommets sont biffés}

\ill{} $0$-connexe \qquad $W_{2}$

\begin{tikzcd}
a \arrow[r, bend left=14] \arrow[r, bend right=14] & b
\end{tikzcd}

Soit \struck{$W$ \ill{} plus fine que tt \ill{}} \quad Soit $W_{2}$ \ill{}

$X$ $0$-connexe, mais pas \struck{filtrant} \ill{}
\[
  \Longrightarrow\ \exists\, X \longrightarrow X' \text{ cofinal},
\]
avec $X' \sim_{W_{\infty}} S^{2}$, ou du moins avec
\[
  H^{2}(X') \simeq H^{2}(S^{2}) ;
\]
donc si $W \subset W_{2}$, plus gén. si \struck{\ill{}} $f\in W \Rightarrow$
\[
  f^{*} : H^{2}(T,\underline{k}) \xrightarrow{\ \sim\ } H^{2}(S,\underline{k})
\]
($k$ un \ill{} groupe \uncertain{commutatif}) \ill{} des $X$ \ill{} pas tot.
$W$-asph.
\note{les deux membres du morphisme portent $T$ et $S$ ; l'un et l'autre sont
tracés d'un seul jet et la lecture en est douteuse}

\page{69}
\note{feuillet scanné à l'envers}

\begin{tikzcd}
X \arrow[r, "u\in W"] \arrow[d, "i\in W_{S}''"'] & Y \arrow[d, "j\in W_{S}''"] \\
\bar X \arrow[r, "\bar u\in W"'] & \bar Y
\end{tikzcd}

\note{de $\bar X$ et $\bar Y$ partent deux flèches vers $S$, étiquetées
$p\in\mathcal{G}$ et $q\in\mathcal{G}$}

\begin{tikzcd}
X \arrow[r, "W_{S}''"] \arrow[d] & \mathcal{G}(X) \arrow[ld] \\
S &
\end{tikzcd}
\note{l'étiquette de la flèche horizontale est précédée d'un « $i\in$ »
biffé ; la flèche verticale en porte une, biffée, illisible. À côté du
diagramme : « $/F_{W}$ » et $\bar u \in W_{S}''$}

\marginal{VI \quad 171-205}

\begin{tikzcd}
X' \arrow[r] \arrow[d] & X \arrow[d] \\
S' \arrow[r, "w"] & S
\end{tikzcd}
\note{sur le feuillet les deux flèches horizontales sont tracées vers la
gauche ; je les redresse pour que le carré se lise}

\[
  \xi \qquad \xi' \qquad H^{p}(S',\xi') \longrightarrow H^{p}(S,\xi)
\]
\struck{$H^{p}(S,\xi)$}
\[
  S\times T \longleftarrow S'\times T
\]

\page{70}
\note{il numérote ce feuillet « 6 »}

\begin{tikzcd}
X_{/F} \arrow[r, hook, "i"] & X' \arrow[r] & e
\end{tikzcd}
\note{un $F$ est porté au-dessus du premier $X$}

\[
  X_{/e} \simeq X_{/F}
\]

$F$ asphérique $\Longleftrightarrow$ $i$ \uncertain{cocofinal}.

Si $X^{o}$ est $W$-asph., alors ($X'^{o}$ dans) $X'$ est $W$-asph., dès que
$i\in W$, et si $W$ satisfait Loc(6), \ill{} est $W$-asph. \ill{} $X'_{/e}$
$W$-asph. \ill{} $X_{/F}$ $W$-asph.

Inversement, supposons que tt objet $0$-connexe de $X^{\wedge}$ soit
$W$-asph. \note{un $F$ est porté au-dessus du signe} Soit
\[
  X \longrightarrow Y \quad \text{cocofinal},
\]
i.e. les $X_{/y}$ $0$-connexes. \struck{Il \ill{} que \ill{}}
Mais $X_{/y}$ est \ill{} $0$-connexe, il est $W$-asph. Donc $X\to Y$ est
$W$-asph. dès \ill{} $\in W$, donc $Y$ est $W$-asph. Ainsi :

\textbf{Prop.} Si \struck{$W$ \ill{}} Pour que $X^{o}$ soit \struck{tot.}
$W$-asphérique il \ill{} et s. que tt objet $0$-connexe de $X^{\wedge}$ soit
$W$-asph., et elle est aussi nécessaire si \struck{\ill{}} \ill{} Loc(6).

\textbf{Cor.} (Plus \struck{ou moins} nécessaire) Si \struck{$X^{o}$ est
cofiltrant et} $W$ satisfait Loc(8), alors pour \ill{} $X$ telle que $X^{o}$
soit filtrant (ou $X$ cofiltrant), les objets $0$-connexes de $X^{\wedge}$
sont (dès que $X_{/F}$ soit \uncertain{cofiltrant dans} \ill{}) $W$-asphériques,
donc $X$ est tot. $W$-asphérique.

\page{71}

\note{le losange qui occupe le haut du feuillet est écrit tête-bêche par
rapport au reste ; c'est, aux étiquettes près, celui du feuillet 79}

\begin{tikzcd}[column sep=small, row sep=small, nodes={font=\scriptsize}]
& T \arrow[ld] \arrow[rd] & \\
\beta/\!/\tilde X \arrow[rd, "\text{cocart.}"'] & & \gamma/\!/\tilde X \arrow[ld] \\
& \tilde X \arrow[ld, "X\ \text{asph.}"'] \arrow[rd] & \\
X \arrow[rd, "W\text{-cocart.}"'] & & e \arrow[ld] \\
& X' &
\end{tikzcd}

\[
  \tilde X \sim \Sigma(T), \quad\text{si } \beta/\!/\tilde X,\ \gamma/\!/\tilde X
  \ \text{asph. (à vérifier)}
\]
\[
  X' \sim \Sigma(\tilde X)
\]
\[
  (vu_{1},\alpha) \sim (vu_{2},\beta)
\]
\[
  (vu_{1},\alpha) \sim (vu_{2},\beta)
\]
\note{les deux lignes sont écrites l'une sous l'autre et, autant que je
puisse voir, identiques ; la seconde porte peut-être $u_{2}$ et $u_{1}$
permutés}

\begin{tikzcd}[column sep=small, row sep=small]
\alpha \arrow[r, "u_{1}'"] \arrow[rd, "u_{1}''"] & B \\
\beta \arrow[ru, "u_{2}'"] \arrow[r, "u_{2}''"'] & C
\end{tikzcd}

\begin{tikzcd}[column sep=small, row sep=small]
a \arrow[r, "u_{1}"] \arrow[rd, "w"'] & b \arrow[d] \\
& x
\end{tikzcd}

\[
  (v,\alpha) \sim (w,\beta)
\]
\ill{} que $v$ \ill{} se factorise par $u_{1}$, \ill{} $w$ \ill{}

\note{plus bas, trois flèches isolées : $\alpha\to\xi$, $\beta\to\xi$,
$B\to\xi$, $C\to\xi$, et la lettre $E$}

\page{72}
\note{il numérote ce feuillet « 7 »}

On voudrait une réciproque, dans des conditions convenables.

\textbf{Lemme.} $X$ \struck{\ill{}}, $F\in X^{\wedge}$, alors $X_{/F}$ irréd.
\note{sous l'énoncé, entre accolades, la définition : « $X_{/F}$ non vide et
\ill{} »}

\begin{tikzcd}[column sep=small, row sep=small]
a \arrow[rd] & \\
& F \\
b \arrow[ru] &
\end{tikzcd}

$a$, $b \in X$, on a $a\times_{F}b \neq \emptyset$, i.e. \ill{}

\begin{tikzcd}[column sep=small, row sep=small]
& a \arrow[rd] & \\
c \arrow[ru, dashed] \arrow[rd, dashed] & & F \\
& b \arrow[ru] &
\end{tikzcd}

\ill{} $c\in\mathrm{Ob}\,X$, rendant le diagramme commutatif.

\textbf{Corollaire.} Pour que $X$ soit \struck{loc.} \struck{\ill{}}
cofiltrant, (i.e. conv. comm. cofiltrante) \struck{\ill{} et \ill{}} \ill{}
et s. que pour \ill{} dans $X$,

\begin{tikzcd}[column sep=small, row sep=small]
& \xi & \\
a \arrow[ru, "\lambda"] \arrow[rd, "\lambda'"'] & & b \arrow[lu, "\mu"'] \arrow[ld, "\mu'"] \\
& \xi' &
\end{tikzcd}

\note{sous le losange, deux flèches pointillées descendent vers un sommet $x$}

\textbf{N.B.} \ill{} on peut compléter en un diagramme commutatif.
\note{la remarque est portée en biais, en travers du coin inférieur gauche ;
je n'en lis que la fin}

Soit $S$ la catégorie ordonnée
\[
  S = \bigl(\, a \to \xi \to b \ ;\ a \to \xi' \to b \,\bigr)
\]

\page{73}

\[
  B = (u_{1},\alpha) \sim (u_{2},\beta)
\]
\[
  C = (u_{2},\alpha) \sim (u_{1},\beta)
\]
\[
  B \longrightarrow \xi
\]

\begin{tikzcd}
\alpha \arrow[r, bend left=16, "u_{1}"] \arrow[r, bend right=16, "u_{2}"'] & b \arrow[r, "v"] & x
\end{tikzcd}

\[
  \xi = (vu_{1},\alpha) = (vu_{2},\beta)
\]
\[
  \xi = (vu_{2},\alpha) = (vu_{1},\beta)
\]

\begin{tikzcd}
a \arrow[r, bend left=16, "v"] \arrow[r, bend right=16, "w"'] & x
\end{tikzcd}

\[
  \xi = (v,\alpha) = (w,\beta)
\]

donc $\mathrm{Hom}(b,x) \neq \emptyset$.

\begin{tikzcd}[column sep=small, row sep=small]
\alpha \arrow[r, "\tilde u"] & \tilde\xi \\
\beta \arrow[r, "\tilde w"'] & \tilde\xi
\end{tikzcd}

il \ill{} \ill{} $v$ se factorise par \ill{} se factorise par $u_{1}$, \ill{}
par $u_{2}$,

\page{74}
\note{il numérote ce feuillet « 8 »}

$Z$ la catégorie obtenue en ajoutant un objet final (pt.\ final) $e$, d'où

\begin{tikzcd}
S \arrow[r, hook, "\text{cocof.}"] \arrow[d] & Z \arrow[r, "\text{ou}"] \arrow[d] & e \arrow[d, no head] \\
X \arrow[r, hook] & X' \arrow[r, "\text{ou}"'] & e
\end{tikzcd}
\note{il étiquette ce carré (D) dans la marge de gauche ; le trait sans tête
de la dernière colonne est, sur le feuillet, une double barre — l'identité}

avec
\[
  \mathrm{Hom}_{X'\supset X}(x,e) \simeq \pi_{0}\bigl({}_{x}\backslash S_{/e}\bigr)
  \simeq \pi_{0}\bigl({}_{x}\backslash S\bigr)
\]
\[
  \simeq \varinjlim \Bigl(\mathrm{Hom}(x,a) \to \mathrm{Hom}(x,\xi)
  \leftarrow \mathrm{Hom}(x,b) \to \mathrm{Hom}(x,\xi') \leftarrow \Bigr)
\]
\note{le colimite est écrite comme un losange de quatre flèches, de sommets
$\mathrm{Hom}(x,a)$, $\mathrm{Hom}(x,\xi)$, $\mathrm{Hom}(x,\xi')$ et
$\mathrm{Hom}(x,b)$, les deux premières partant de $\mathrm{Hom}(x,a)$ et les
deux dernières arrivant de $\mathrm{Hom}(x,b)$}
\[
  \simeq \underbrace{\bigl(\mathrm{Hom}(x,\xi)\amalg\mathrm{Hom}(x,\xi')\bigr)}_{E(x)}
  \big/ R_{(x)},
\]
$R_{(x)}$ la relation d'équivalence engendrée par les deux relations
\[
  (*) \quad
  \begin{cases}
    \lambda v \sim \lambda' v & v : x \to a \\
    \mu w \sim \mu' w & w : x \to b
  \end{cases}
\]

Le foncteur $S \to X$ se \uncertain{prolonge} canoniquement en
$S \to X_{/e}$ \ill{} (grâce au fait que $e$ est un objet final de $Z$, cf.
diagramme (D)), soit \ill{}

\page{75}
\note{feuillet de brouillon : deux blocs y sont écrits tête-bêche l'un par
rapport à l'autre, et plusieurs sont raturés}

\begin{tikzcd}
\alpha \arrow[r, bend left=16] \arrow[r, bend right=16] & \beta \arrow[r] & \gamma
\end{tikzcd}

\begin{tikzcd}
a \arrow[r, bend left=16] \arrow[r, bend right=16] & b \arrow[r] & c
\end{tikzcd}
\note{sous ce second diagramme, une flèche montante issue d'un sommet $a$}

\begin{tikzcd}
\alpha \arrow[r, bend left=16] \arrow[r, bend right=16] & \beta \arrow[r] & C
\end{tikzcd}
\note{ce troisième diagramme est entièrement raturé}

\[
  \mathrm{Hom}(b,\xi) \rightrightarrows \mathrm{Hom}(a,\cdot)
\]

\begin{tikzcd}
X_{/\!/b} & \tilde X \arrow[l] & X_{/\!/a} \arrow[l, "\text{cocart.}"']
\end{tikzcd}
\note{les deux indices sont tracés d'un jambage et restent douteux}

\note{on lit encore, isolés et pour partie biffés, $\tilde X$, $X_{/\!/e}$,
$e$ et une flèche $e \to X$}

\page{76}
\note{il numérote ce feuillet « 9 »}

\begin{tikzcd}[column sep=small, row sep=small]
& \tilde\xi & \\
\tilde a \arrow[ru, "\tilde\lambda"] \arrow[rd, "\tilde\lambda'"'] & & \tilde b \arrow[lu, "\tilde\mu"'] \arrow[ld, "\tilde\mu'"] \\
& \tilde\xi' &
\end{tikzcd}

le diagramme correspondant dans $\tilde X = X_{/e}$.

\textbf{Lemme.} On a \struck{\ill{}}
\[
  \tilde X = \tilde X_{/\!/\tilde\xi} \ \cup\ \tilde X_{/\!/\tilde\xi'}
\]
\[
  \tilde X_{/\!/\tilde\xi} \ \cap\ \tilde X_{/\!/\tilde\xi'}
  = \tilde X_{/\!/\tilde a} \ \cup\ \tilde X_{/\!/\tilde b}
\]
\note{le second indice des deux premiers membres est tracé d'un jambage qui
ne se distingue guère de celui de $\tilde\xi$ ; c'est le losange ci-dessus qui
fixe la lecture $\tilde\xi'$}

La première relation est \uncertain{claire}. Pour la deuxième, \ill{} soit

\begin{tikzcd}[column sep=small, row sep=small]
\tilde a \arrow[r, "\tilde u"] & \tilde\xi \\
\tilde u' \arrow[r] & \tilde\xi'
\end{tikzcd}

dans $\tilde X$, au-dessus de

\begin{tikzcd}[column sep=small, row sep=small]
& \xi \\
a \arrow[ru, "u"] \arrow[rd, "u'"'] & \\
& y
\end{tikzcd}

dans $X$.
\note{les deux petits diagrammes sont tracés en diagonale ; les sources des
flèches inférieures ne se lisent pas avec certitude}

\page{77}
\note{ce feuillet reprend, avec la catégorie plus simple $D = (a
\rightrightarrows b)$, l'argument que le feuillet 74 mène avec $S$ ; c'est une
mouture antérieure, non une répétition. Un « 6 » cerclé, barré d'un trait,
figure en tête}

\note{deux accolades encadrent la ligne $e \to a \rightrightarrows b$ :
$Z$ au-dessus, $D$ au-dessous}

\begin{tikzcd}
D \arrow[r, "\text{cofinal}"] \arrow[d] & Z \arrow[r, "\text{ou}"] \arrow[d] & e \arrow[d, no head] \\
X \arrow[r, hook, "\text{cofinal}"'] & X' \arrow[r, "\text{ou}"'] & e
\end{tikzcd}
\note{sous l'étiquette « cofinal » de la flèche supérieure court un second
mot, biffé, que je ne lis pas ; le trait sans tête de la troisième colonne
est, sur le feuillet, une double barre}

\struck{$e\backslash X$ \ill{}}

\begin{tikzcd}
a \arrow[r, bend left=14] \arrow[r, bend right=14] & b
\end{tikzcd}

\struck{\ill{}}

\[
  \mathrm{Hom}(e,x) \simeq \pi_{0}\bigl({}_{e}\backslash D_{/x}\bigr)
  \simeq \pi_{0}\bigl(D_{/x}\bigr)
  \simeq \varinjlim \bigl(\mathrm{Hom}(b,x) \rightrightarrows \mathrm{Hom}(a,x)\bigr)
\]

un objet de ${}_{e}\backslash X_{/e}$ est déterminé par un objet \uncertain{$x$}
de $X$, et une flèche $a \to x$. L'objet de ${}_{e}\backslash X_{/e}$ défini
par \ill{} \marginal{id} \ill{} s'envoie dans tous les autres, donc
${}_{e}\backslash X_{/e}$ est un \ill{}, donc $X \to X'$ est cofinal.

\page{78}
\note{il numérote ce feuillet « 10 »}

Car donc
\[
  u \in \mathrm{Hom}(x,\xi) \subset \mathrm{Hom}(x,\xi)\amalg\mathrm{Hom}(x,\xi')
  \overset{\text{déf}}{=} E(x)
\]
\[
  u' \in \mathrm{Hom}(x,\xi') \subset E(x)
\]
et $u \sim u' \bmod R(x)$. \quad (\ill{})

Donc \struck{il existe} comme $u \neq u'$, il existe une chaîne d'objets de
$E(x)$
\[
  f_{0} = u,\ f_{1},\ f_{2},\ \ldots\ f_{n} = u'
\]
($n \geq 1$), telle que pour $0 \leq i \leq n-1$, \struck{$f_{i}$} $f_{i}$ et
$f_{i+1}$ soient reliés \ill{} par une des deux relations $(*)$. En
particulier, \struck{\ill{}} pour $i = 0$, il existe $v : x \to a$ tel que
$f_{0} = \lambda v$ \struck{et} (et $f_{1} = \lambda' v$), ou bien
$w : x \to b$ tel que $f_{0} = \mu w$ (et $f_{1} = \mu' w$).

Donc $u$ se factorise soit par $(a,\lambda)$, soit par $(b,\mu)$. Mais cela
signifie \ill{} que $\tilde u : \tilde x \to \tilde\xi$ se factorise soit par
$\tilde a$, soit par $\tilde b$, \struck{\ill{}} \struck{$\tilde\xi$}
\uncertain{$\tilde\xi$} est dans $\tilde X_{/\!/\tilde a} \cup
\tilde X_{/\!/\tilde b}$, qed.
\note{le sujet de la dernière proposition devrait être $\tilde x$ et non
$\tilde\xi$ ; je transcris ce qui est tracé}

\note{en marge gauche, le losange $\tilde a \to \tilde\xi \leftarrow \tilde b$,
$\tilde a \to \tilde\xi' \leftarrow \tilde b$, avec les étiquettes
$\lambda$, $\mu$, $\lambda'$, $\mu'$, et au-dessous $u$, $u'$, $x$}

\page{79}
\note{il numérote ce feuillet « 11 »}

\textbf{Cor.} $\tilde X$ est $0$-connexe, donc (si les objets $0$-connexes de
$X^{\wedge}$ sont $W$-asph., pour un $W$ donné) il est $W$-asphérique.

(si $X \hookrightarrow X'$ est cofinal).

\begin{tikzcd}[column sep=small, row sep=small, nodes={font=\scriptsize}]
& T = \tilde X_{/\!/(\tilde a,\tilde b)} \arrow[ld] \arrow[rd] & \\
\tilde X_{/\!/\tilde b} \arrow[rd, "W\text{-cocart.}"'] & & \tilde X_{/\!/\tilde a} \arrow[ld] \\
& \simeq \Sigma^{1}(T) \arrow[ld] \arrow[rd] & \\
\tilde X_{/\!/\tilde\xi} \arrow[rd, "W\text{-cocart.}"'] & & \tilde X_{/\!/\tilde\xi'} \arrow[ld] \\
& \tilde X \simeq \Sigma^{1}_{W}(T) \arrow[ld] \arrow[rd] & \\
X \arrow[rd] & & e \arrow[ld] \\
& X' &
\end{tikzcd}
\note{le sommet du milieu porte, avant le signe $\simeq$, un membre
entièrement biffé}
\note{en marge des sommets : « $0$-connexes donc $W$-asph. » pour
$\tilde X_{/\!/\tilde\xi'}$, « $0$-connexe donc $W$-asph. » pour $\tilde X$,
« $0$-connexe donc $W$-asph. par \ill{} » pour $X$ ; et, à droite de
$\tilde X$, « (\ill{} Loc(6a)) »}

\[
  X' \simeq \Sigma^{1}_{W}(X) \simeq \Sigma^{2}_{W}(\tilde X)
  \simeq \Sigma^{3}(T)
\]
\note{l'argument du terme médian est tracé d'un seul jambage : je le lis
$\tilde X$, ce que la chaîne $\tilde X \simeq \Sigma^{1}(T)$,
$X \simeq \Sigma^{1}(\tilde X)$ du feuillet 66 rend seule cohérente, mais la
lettre reste douteuse}

Si $T = \emptyset$, alors $X' \sim_{W_{\infty}} \Sigma^{2}(E) = S^{2}$, donc
si $W$ plus fine que $H^{2}(-,k)$, $k$ \uncertain{groupe commutatif} un \ill{},
alors $X'$ n'est pas $W$-asph., \ill{}. Donc dans ce cas on \ill{} doit avoir
\[
  T \neq \emptyset .
\]

\end{document}
